% Part: history
% Chapter: biographies 
% Section: georg-cantor

\documentclass[../../../include/open-logic-section]{subfiles}

\begin{document}

\olfileid{his}{bio}{can}

\olsection{ګېورګ کانتور}

\olphoto{cantor-georg}{ګېورګ کانتور}

د ګېورګ کانتور، چې نوم ئې (\textsc{ګې}-اورګ
\textsc{کان}-تور) ويل کېږي، يوه پخواني ژوندليک
ادعا کړې وه چې هغه په يوه بېړۍ کښې زېږېدلے او
موندل شوے و، چې د روسيې سېنټ پېټرزبرګ ته روانه
وه، او مور او پلار ئې نۀ وو معلوم۔ خو دا رښتيا
نۀ دي؛ که څه هم هغه په 1845 کښې په سېنټ پېټرزبرګ
کښې زېږېدلے و۔

کانتور په 1867 کښې د برلين له پوهنتون څخه په
رياضياتو کښې دوکتورا ترلاسه کړه۔ هغه د سټونو
په تيورۍ کښې د خپل کار له امله پېژندل کېږي،
او د سټونو د تيورۍ د يوې بېلې څېړنيزې څانګې
په توګه بنسټ اېښودل ورته منسوبېږي۔ هغه لومړنے
کس و چې ثابت ئې کړل د بېلابېلو اندازو
نامتناهي سټونه شته۔ د هغه تيوريو، او په تېره
بيا د نامتناهياتو تيورۍ ئې، د هغه وخت د
رياضيپوهانو تر منځ ډېر بحث پيدا کړ، او کار
ئې د اختلاف وړ و۔

د کانتور مذهبي عقيدې او رياضيکي کار يو له بل
سره نۀ بېلېدونکي تړلي وو؛ آن هغه دا ادعا هم
کوله چې د ترانسفينيټ عددونو تيوري خداے نېغ
په نېغه ورته رسولې وه۔ د ژوند په وروستيو
کلونو کښې کانتور په ذهني ناروغۍ اخته و۔
له 1894 څخه پيل او په وروستيو کلونو کښې
ئې په زيات تکرار، کانتور روغتون ته داخلېدو۔
پر کار ئې سختو نيوکو، چې له رياضيپوه
لېوپولډ کرونېکر سره اختلاف هم پکښې و،
هغه ژور خپګان او له رياضياتو سره د شوق
نشتوالي ته ورساوۀ۔ د ژور خپګان په دورو
کښې کانتور فلسفې او ادب ته مخه کوله، او
آن يوه تيوري ئې هم خپره کړه چې فرانسس
بېکن د شېکسپير د ډرامو ليکوال و۔

کانتور د 1918 د جنورۍ پر 6مه په هالې
کښې په يوه درملنځاے کښې وفات شو۔

\begin{reading}
د کانتور د بشپړو ژوندليکونو دپاره
\citet{Dauben1990} او \citet{Grattan-Guinness1971}
وګورئ۔ د کانتور انقلابي نظرونه د بي بي سي
راډيو 4 په پروګرام \emph{د رياضياتو لنډ تاريخ}
کښې هم بيان شوي دي \citep{Sautoy2014}۔ که
غواړئ د کانتور د تيوريو په هکله د رېپ
په بڼه واورئ، \citet{Rose2012} وګورئ۔
\end{reading}

\end{document}
